%=======================================================================
%  READER'S NOTES -- the environment, and the registry that keeps their
%  text out of the chapter bodies.
%
%  A reader's note is a long-form expansion of a step the running text
%  compresses: the algebra written out with nothing assumed.  The notes
%  are set in from the margin, in a smaller face, behind a coloured rule,
%  so that the spine of the argument can be read straight through and the
%  arithmetic is there for whoever wants it.
%
%  The notes themselves live in readernotes_ch<N>.tex, one file per
%  chapter.  A chapter body calls
%
%       \rnote{ch6:6.131-combining}
%
%  at the point where the note belongs, and the matching
%
%       \rnotedef{ch6:6.131-combining}{%
%       \begin{readernote}[the five substitutions ...]
%       ...
%       \end{readernote}%
%       }
%
%  in the notes file supplies the text.  Moving a note is therefore a
%  matter of moving one line, and rewriting one never touches the
%  chapter.  A \rnote whose key was never registered prints a visible
%  red placeholder and warns, rather than vanishing silently.
%
%  Requires: xcolor.  Input from the preamble, before the notes files.
%=======================================================================

\makeatletter

% ------------------------------- the environment -----------------------
% The optional argument is the note's title, shown in parentheses after
% the words "Reader's note".  It may be omitted.
% ------------------------------- the two inks --------------------------
% The whole note -- prose, algebra, boxes and rules alike -- is set in
% rnink, so that a note reads as one object and the chapter's own text is
% the only black on the page.  The reason column is set in rnfaint, the
% same ink lightened, which keeps the algebra ahead of its justification
% without introducing a second hue.  Change these two lines to recolour
% every note in the book.
\colorlet{rnink}{blue!45!black}
\colorlet{rnfaint}{blue!45!black!58!white}

\newenvironment{readernote}[1][]{%
  \par\smallskip
  \def\rn@arg{#1}\def\rn@empty{}%
  \begingroup
  \advance\leftskip by 1.2em
  \tolerance=1400 \emergencystretch=3em
  % Tighten the gap between the algebra and the reason column.  NOTE:
  % amsmath's \minalignsep is a *macro* expanding to 10pt, not a dimen
  % register, so `\minalignsep=6pt' would typeset the words `10pt=6pt'
  % into the page rather than assigning anything.  It must be redefined.
  \renewcommand{\minalignsep}{6pt}%
  \small\color{rnink}\noindent
  \rule[-0.18em]{1.6pt}{0.95em}\enspace
  {\scshape Reader's note}%
  \ifx\rn@arg\rn@empty\else\ \textit{(#1)}\fi
  {\bfseries.}\enspace\ignorespaces
}{\par\endgroup\smallskip}

% --------------------------------- the registry ------------------------
% \rnotedef{key}{text} stores the text; \newcommand makes \rnotedef long,
% so blank lines inside the text are allowed.  The stored macro takes no
% arguments, so no \long is needed on the \gdef itself.
\newcommand{\rnotedef}[2]{%
  \@ifundefined{rn@text@#1}{}%
    {\@latex@warning@no@line{Reader's note `#1' registered twice}}%
  \expandafter\gdef\csname rn@text@#1\endcsname{#2}}

\newcommand{\rnote}[1]{%
  \@ifundefined{rn@text@#1}%
    {\@latex@warning@no@line{Reader's note `#1' is undefined}%
     \par\smallskip\noindent
     \textcolor{red}{[reader's note \texttt{#1} missing]}\par\smallskip}%
    {\csname rn@text@#1\endcsname}}

% ------------------------------ notebook furniture ---------------------
% The notes are written as a working notebook is written: a marker, a
% ledger of equalities with the reason for each in the right-hand column,
% a conclusion.  Prose appears only where a line of algebra cannot carry
% the point.  \rnlab is the marker; it opens a paragraph.
%
%       \rnlab{goal}  ...one line...
%
%       \[\begin{aligned}
%          x &= y   &&\eqr{2.50}\\
%          y &= z   &&\text{product rule}
%       \end{aligned}\]
%
% The reason column is the second alignment pair of amsmath's aligned,
% opened by the literal `&&'.  Write those two ampersands out: an
% alignment tab arriving from a macro is not seen by the alignment
% scanner, so \rnby-style shorthands do not work here.  What follows the
% tabs is wrapped in \rnwhy, which sets the justification a size down and
% in grey, so that the eye reads the algebra first and the reason second:
%
%       x &= y   &&\rnwhy{\eqr{2.50}}\\
%       y &= z   &&\rnwhy{\text{product rule}}
%
% The argument of \rnwhy is ordinary MATH, written exactly as it would be
% in the display itself, with \text{...} round any prose: the macro boxes
% it at \footnotesize and re-opens math inside the box, so nothing about
% the content has to change.  Being an \mbox it will not line-break, which
% is the right behaviour for a marginal reason and a warning when one has
% grown too long to be a reason.
%
% Conclusions use \therefore and \boxed, both already available from
% amssymb and amsmath.
\newcommand{\rnlab}[1]{%
  {\footnotesize\scshape #1}\enspace\ignorespaces}
\newcommand{\rnwhy}[1]{%
  \mbox{\footnotesize\textcolor{rnfaint}{$#1$}}}

\makeatother
